\chapter{ВВЕДЕНИЕ}

Искусственные нейронные сети являются основным инструментом решения задач
классификации, для которых не удаётся явно сформулировать решающее правило. Простейшей
и исторически первой архитектурой такого рода является многослойный персептрон ---
последовательность полносвязных слоёв, разделённых нелинейными функциями активации.
Несмотря на простоту устройства, именно на нём отрабатываются все базовые механизмы
обучения: прямое распространение сигнала, вычисление функции потерь, обратное
распространение ошибки и градиентный спуск.

Практическая часть работы выполнена на языке C++ и включает библиотеку с реализацией
перечисленных механизмов и
серию вычислительных экспериментов по выбору архитектуры сети.

\textbf{Постановка задачи.} Разработать многослойный персептрон с настраиваемой
архитектурой, обучаемый методом обратного распространения ошибки и градиентного спуска,
и применить его к задаче бинарной классификации диагноза по тридцати признакам
 из набора.
